\documentclass[11pt,a4paper]{article}
\usepackage{authblk}
\usepackage[utf8]{inputenc}
\usepackage{amsmath,amssymb,amsthm}
\usepackage{geometry,graphicx,float}
\usepackage{tikz}
\usetikzlibrary{arrows.meta,positioning}
\usepackage{longtable}
\usepackage{multicol}
\geometry{margin=0.9in}
\usepackage{hyperref}
\usepackage{orcidlink}    % For ORCID icons

\usepackage{xcolor}    % always load before hyperref
\usepackage{hyperref}
\usepackage{listings}
\usepackage{booktabs}

% Choose your colors (any named color or custom RGB/hex)
\hypersetup{
    colorlinks=true,           % false = framed boxes, true = colored text
    linkcolor=cyan,            % internal links (ToC, \ref, \pageref)
    citecolor=cyan,           % bibliography citations
    urlcolor=cyan,          % URLs and \url{}
    filecolor=cyan,            % file links
    menucolor=cyan,             % PDF menu links (rarely used)
    % Optional: make links bold or italic
    % linkbordercolor={0 1 0}, % only if colorlinks=false
}

% Or use custom colors:
\definecolor{myblue}{RGB}{0,80,180}
\definecolor{mygreen}{RGB}{0,120,0}
\hypersetup{
    colorlinks=true,
    linkcolor=cyan,
    citecolor=cyan,
    urlcolor=cyan,
}

%\newline email: ntelis.pierros@gmail.com

%\author[$\dagger$]{Anonymous}
%\affil[$\dagger$]{Anonymous Affiliation}

\renewcommand\Affilfont{\itshape\small}

\date{\today}% \\ Conceived: 08 November 2025}

% === FULLY TRANSPARENT THEOREM ENVIRONMENTS ===
\newtheorem{primitive}{Primitive}
\newtheorem{axiom}{Axiom}
\newtheorem{construct}{Construct}
\newtheorem{definition}{Definition}
\newtheorem{theorem}{Theorem}
\newtheorem{lemma}{Lemma}
%\newtheorem{proof}{Proof}
\newtheorem{remark}{Remark}


\author[$\dagger$,$\ddagger$]{Pierros Ntelis\orcidlink{0000-0002-7849-2418}}
%\affil[$\dagger$]{Independent Researcher, 13 Allée Turcat Méry, 13008 Marseille, France}
\affil[$\dagger$]{School of Physics, Harbin Institute of Technology, Harbin 150001, People's Republic of China}
\affil[$\ddagger$]{Institute of Theoretical Physics, National University of Uzbekistan, Tashkent 100174, Uzbekistan}





\lstset{
    basicstyle=\ttfamily\small,
    backgroundcolor=\color{gray!10},
    frame=single,
    numbers=left,
    numberstyle=\tiny,
    breaklines=true,
    captionpos=b,
}

\title{A temporal-spatial-entropic metric and its Kretschmann scalar}
\date{\today; conceived: 5 June 2026 }



\begin{document}

\maketitle

\begin{abstract}
We present \href{https://github.com/lontelis/Einsteinpy_to_latex}{{\color{cyan}\texttt{EinsteinKPy2Tex}}}, 
an open-source Python package that 
computes fundamental geometric quantities for general relativistic metrics. 
Built upon the EinsteinPy library, the code calculates Christoffel symbols, 
Riemann curvature tensor, Ricci tensor and scalar, Einstein tensor, geodesic 
equations, and the Kretschmann scalar. Results are exported as professionally 
formatted \LaTeX\ documents with only non-zero components, facilitating 
direct inclusion in research papers and educational materials.

The software includes pre-configured implementations for two cornerstone 
metrics: the Friedmann-Lema\^{i}tre-Robertson-Walker (FLRW) metric for 
cosmology and the Schwarzschild metric for black hole physics. Built-in 
citations for Einstein (1915) and the EinsteinPy library (Ribeiro et al. 
2020) ensure proper academic attribution.

Building upon the Advanced Manifold-Metric Pairs (AMMP) framework recently 
developed by Ntelis (2025), we introduce and analyze a novel 5-dimensional 
metric that promotes entropy to a genuine spacetime coordinate—the txyzS 
metric. This represents the first concrete application of the AMMP 
framework, demonstrating that thermodynamic quantities can be treated as 
coordinates on equal footing with space and time. For this new metric, we 
compute all geometric quantities and derive the Kretschmann scalar, which 
acquires additional terms $4(\beta S)^4$ and $12H^2(\beta S)^2$ arising 
from the entropy dimension.

Numerical simulations of the Kretschmann scalar for matter-dominated, 
radiation-dominated, and $\Lambda$-dominated cosmologies reveal that the 
entropy dimension amplifies spacetime curvature by a factor ranging from 
$1.01$ to $4.15$ for entropy coordinates $S = 0.5$ to $2.0$, with the largest 
deviations occurring at late times ($t \sim 1$). The ratio 
$K_{\text{txyzS}}/K_{\text{FLRW}}$ exceeds unity for all $S > 0$, providing 
a clear observational target for constraining the AMMP framework using 
cosmic microwave background and large-scale structure data.

\href{https://github.com/lontelis/Einsteinpy_to_latex}{{\color{cyan}\texttt{EinsteinKPy2Tex}}} 
is designed for educators, students, and 
researchers needing rapid, accurate tensor calculations with publication-ready 
output. The entropy-dimension implementation establishes a new paradigm 
for incorporating thermodynamic degrees of freedom into general relativity, 
with potential applications to black hole thermodynamics, quantum gravity, 
and early-universe cosmology. The software is open-source under an MIT 
license and permanently archived on Zenodo (DOI: 10.5281/zenodo.20563701), 
with the latest version available at \href{https://github.com/lontelis/Einsteinpy_to_latex}{{\color{cyan}\texttt{EinsteinKPy2Tex}}}.
\end{abstract}

\iffalse
\begin{abstract}
We present \href{https://github.com/lontelis/Einsteinpy_to_latex}{{\color{cyan}\texttt{EinsteinKPy2Tex}}}, an open-source Python package that computes fundamental geometric quantities for the Friedmann-Lema\^{i}tre-Robertson-Walker (FLRW) and Schwarzschild metrics using the EinsteinPy library. The code automatically calculates Christoffel symbols, Riemann curvature tensor, Ricci tensor, Ricci scalar, Einstein tensor, geodesic equations, and Kretschmann scalar. Results are exported as professionally formatted \LaTeX\ documents, facilitating direct inclusion in research papers and educational materials.

The FLRW implementation is essential for cosmology, enabling the study of the universe's large-scale structure and dynamics. The Schwarzschild implementation is fundamental for black hole physics, describing the spacetime geometry around non-rotating compact objects. The code includes built-in citations for the original Einstein paper (1915) and the EinsteinPy library (Ribeiro et al. 2020), ensuring proper academic attribution.

\href{https://github.com/lontelis/Einsteinpy_to_latex}{{\color{cyan}\texttt{EinsteinKPy2Tex}}} is designed for educators, students, and researchers needing rapid, accurate tensor calculations with publication-ready output. The software is available at \href{https://github.com/lontelis/Einsteinpy_to_latex}{{\color{cyan}\texttt{GitHub}}} under an MIT license.
\end{abstract}
\fi

\keywords{ \textbf{Keywords:} General Relativity ; Schwarzschild metric ; FLRW metric ;
          Kretschmann scalar ; Geodesic equations ; Tensor calculus ;
          \LaTeX\ export ; Symbolic computation ; Advanced Manifold-Metric Pairs;  
          Entropy; Cosmology ; Black hole physics; EinsteinPy ; EinsteinKPy2Tex}

\tableofcontents

\section{Introduction}

General relativity (GR) remains the cornerstone of modern astrophysics, from cosmological models of the expanding Universe to the description of black hole spacetimes. However, calculating the fundamental tensors---Christoffel symbols, the Riemann curvature tensor, the Ricci tensor and scalar, and the Einstein tensor---for even moderately complex metrics is error-prone and time-consuming when performed manually. While symbolic computation packages like SymPy \cite{Meurer2017} and specialized libraries like EinsteinPy \cite{Ribeiro2020} exist, there remains a gap between symbolic calculation and publication-ready output. Researchers often need to verify intermediate steps, derive geodesic equations, or compute curvature invariants such as the Kretschmann scalar, all while maintaining a clear record of their calculations for papers and theses.

This need is particularly acute in cosmology and black hole astrophysics, where sophisticated software infrastructure has become essential for modern research. In cosmology, the community has developed numerous public codes for computing cosmic microwave background (CMB) anisotropies and large-scale structure. Notable examples include \texttt{CAMB} (Code for Anisotropies in the Microwave Background) \cite{Lewis2000}, \texttt{CLASS} (Cosmic Linear Anisotropy Solving System) \cite{Blas2011, Lesgourgues2011a, Lesgourgues2011b}, and the earlier \texttt{CMBFAST} \cite{Seljak1996}. These Boltzmann codes solve the full Einstein-Boltzmann hierarchy and have enabled precision cosmology, from the WMAP era to Planck \cite{Planck2018}. They are often used in conjunction with parameter estimation frameworks such as \texttt{CosmoMC} \cite{Lewis2002} and \texttt{MontePython} \cite{Audren2013}, and more recently with neural network-based emulators like \texttt{CosmoPower} \cite{SpurioMancini2021} that accelerate cosmological inference.

In parallel, the black hole and gravitational wave communities have developed their own specialized software for computing geodesics, orbital trajectories, and gravitational waveforms. While these existing codes are powerful, they are typically designed for large-scale numerical computations, not for educational purposes or for generating clean, publication-ready symbolic output.

\href{https://github.com/lontelis/Einsteinpy_to_latex}{{\color{cyan}\texttt{EinsteinKPy2Tex}}} fills this niche. It is a lightweight Python tool that bridges the gap between symbolic tensor calculations and \LaTeX\ export, specifically designed for the FLRW and Schwarzschild metrics. The code:
\begin{itemize}
    \item Computes all standard GR tensors (metric, Christoffel, Riemann, Ricci, Einstein)
    \item Derives the geodesic equations automatically from the Christoffel symbols
    \item Calculates the Kretschmann scalar $K = R^{\mu\nu\rho\sigma}R_{\mu\nu\rho\sigma}$, a key curvature invariant for identifying singularities
    \item Exports all non-zero components as ready-to-compile \LaTeX\ documents with proper index spacing
    \item Includes built-in citations for Einstein (1915), the EinsteinPy library \cite{Ribeiro2020}, and this software
\end{itemize}

The FLRW implementation enables studies of cosmological dynamics, including the dependence of curvature invariants on the scale factor $a(t)$ and its derivatives. The Schwarzschild implementation provides the foundation for black hole physics, correctly recovering the vacuum solution $R_{\mu\nu}=0$ and the Kretschmann scalar $K = 48 G^2 M^2 / (c^4 r^6)$, which signals the physical singularity at $r=0$.

\href{https://github.com/lontelis/Einsteinpy_to_latex}{{\color{cyan}\texttt{EinsteinKPy2Tex}}} is designed for educators, students, and researchers needing rapid, accurate tensor calculations with publication-ready output. Under an MIT license, the software is available, at \url{https://github.com/lontelis/EinsteinKPy2Tex}.

While the Kretschmann scalar has been computed analytically for the 
Schwarzschild and FLRW metrics since the early days of general relativity 
\cite{Schwarzschild1916, Friedmann1922}, these calculations have 
traditionally been treated as separate textbook exercises. Recent research, 
however, has begun exploring the Kretschmann scalar as a bridge between 
these two regimes. In particular, Dabash \ Big  (2019) \cite{Dabash:2025fpd}
proposed using the Kretschmann scalar to study the Bang–Big-Rip model. 
In our work, we propose to use Kretschmann scalar, to connect local (Schwarzschild) and 
cosmic (FLRW) regimes, deriving geometric constraints. 
Other studies have investigated connections between 
black hole and cosmological singularities, including modeling the FLRW 
universe inside a black hole \cite{Christodoulou2015}. Evidently, 
we explore these novel mathematically, and interesting structures, in the framework of
using curvature invariants to study singularities in modified gravity 
theories \cite{horndeski1974second,kobayashi2011generalized,2012PhR...513....1C,2018FrASS...5...44E,
,Ntelis2023,Ntelis2025_AMMP}. 

This work is grounded in the Advanced Manifold-Metric Pairs (AMMP) 
framework recently introduced by Ntelis (2025) \cite{Ntelis2025_AMMP}. 
The AMMP framework provides 
a mathematical structure for promoting physical quantities—such as entropy, 
temperature, and mass—from mere parameters to genuine spacetime coordinates. 
Unlike traditional extra-dimensional theories (e.g., Kaluza-Klein or string 
theory) which introduce additional spatial dimensions, the AMMP approach 
allows the extra dimensions to represent thermodynamic or information-theoretic 
degrees of freedom. This offers a potential bridge between general relativity 
and thermodynamics, two pillars of modern physics whose deep connection was 
first hinted at by Bekenstein and Hawking but remains incompletely understood.

In this paper, we present the first concrete application of the AMMP framework 
by constructing and analyzing a new metric that includes entropy as an extra 
dimension. Specifically, we extend the standard 4-dimensional FLRW metric 
to 5 dimensions by adding an entropy coordinate $S$, resulting in the 
txyzS metric. We call this metric, an \textit{temporal-spatial-entropic metric}. This choice is motivated by the fundamental role of entropy 
in black hole thermodynamics and cosmology, as well as the mathematical 
tractability of the entropy dimension within the AMMP framework. We compute 
all geometric quantities—Christoffel symbols, Riemann curvature, Ricci 
tensor and scalar, Einstein tensor, geodesic equations, and the Kretschmann 
scalar—for this new metric. We then perform numerical simulations of the 
Kretschmann scalar across different cosmological eras and entropy values, 
comparing the results to standard FLRW predictions. Our analysis reveals 
that the entropy dimension introduces observable deviations from 
conventional cosmology, particularly at early times and for larger entropy 
coordinates. This work thus establishes the AMMP framework as a viable 
approach for generating novel metrics with physical consequences, opening 
new directions for exploring the relationship between gravity and 
thermodynamics.

\href{https://github.com/lontelis/Einsteinpy_to_latex}{{\color{cyan}\texttt{EinsteinKPy2Tex}}} \textit{is the first software package to provide a unified, 
open-source platform for the symbolic computation, numerical verification, 
and direct visual comparison of these invariants.} By automating the 
generation of both symbolic expressions and publication-ready plots, our 
tool lowers the barrier for students and researchers to engage with this 
growing field of research at the interface of black hole and cosmological 
physics. The software enables rapid testing of metric properties, 
visualization of singularity structures, and reproduction of analytical 
results---all with direct \LaTeX\ output for publications.

\section{Software Architecture}

Figure~\ref{fig:flowchart} presents a flowchart illustrating the algorithm implemented in both scripts. 

After starting the code, we first define the important quantities, i.e., the coordinates and symbols. 
We then define the metric components. These definitions are described within the gray area. 

Next, we perform tensor computations via EinsteinPy, followed by the computation of the geodesic equations using symbolic Python code. We then compute the Kretschmann scalar, also using symbolic Python code.
The tensor computations are represented with green colors. 

Finally, using a Python-\LaTeX\ interface, we translate the computations into \LaTeX\ format, represented with pink color.

\begin{figure}[H]
    \centering
    \includegraphics[width=0.85\textwidth]{flowchart.png}
    \caption{Flowchart of \href{https://github.com/lontelis/Einsteinpy_to_latex}{{\color{cyan}\texttt{EinsteinKPy2Tex}}}. The code follows a linear pipeline: metric definition, tensor computation via EinsteinPy, geodesic and Kretschmann calculations, and \LaTeX\ export.}
    \label{fig:flowchart}
\end{figure}


\begin{lstlisting}[caption=High-level algorithm pseudo-code, label=lst:algorithm]
Input: Metric definition (FLRW or Schwarzschild)
Output: LaTeX file with all tensor components

1. Define coordinates: t, x, y, z (FLRW) or t, r, \theta, \phi (Schwarzschild)
2. Define metric parameters: a(t), c (FLRW) or G, M, c, R_s = \frac{2GM}{c^2} (Schwarzschild)
3. Construct metric tensor g_{\mu\nu}
4. Compute using EinsteinPy:
   - Christoffel symbols from metric
   - Riemann tensor from metric
   - Ricci tensor from metric
   - Ricci scalar from metric
   - Einstein tensor from metric
5. Compute geodesic equations: \frac{\dd x^\mu}\dd\lambda^2} + \Gamma^\mu_{\nu\rho} \frac{\dd x^\nu}{\dd \lambda}\frac{\dd x^\rho}{\dd \lambda} = 0
6. Compute Kretschmann scalar: K = R^{\mu\nu\rho\sigma} R_{\mu\nu\rho\sigma}
7. Generate LaTeX output (non-zero components only)
8. Save to .tex file
\end{lstlisting}


\subsection{Dependencies}

The software requires:
\begin{itemize}
    \item Python 3.7 or higher
    \item SymPy \cite{Meurer2017} for symbolic mathematics
    \item NumPy \cite{Harris2020} for numerical array operations
    \item EinsteinPy \cite{Ribeiro2020} for GR tensor operations
\end{itemize}

Installation is straightforward via pip:
\begin{lstlisting}[language=bash]
pip install sympy numpy einsteinpy
\end{lstlisting}

\subsection{Code Structure}

The package consists of two independent scripts:
\begin{itemize}
    \item \texttt{flrw\_EinsteinKPy2Tex.py}: FLRW cosmological metric with scale factor $a(t)$
    \item \texttt{schw\_EinsteinKPy2Tex.py}: Static spherically symmetric metric
\end{itemize}

Each script follows a modular design:
\begin{enumerate}
    \item Metric definition in Cartesian or spherical coordinates
    \item Tensor computation via EinsteinPy API calls
    \item Geodesic equation derivation
    \item Kretschmann scalar calculation
    \item \LaTeX\ output generation with automatic spacing for multi-character indices
\end{enumerate}

\section{Physical Quantities Computed}

The software computes the following quantities for both metrics:

\subsection{Metric Tensor $g_{\mu\nu}$}
The fundamental geometric quantity defining spacetime intervals:
\begin{equation}
ds^2 = g_{\mu\nu} dx^\mu dx^\nu
\end{equation}

\subsection{Christoffel Symbols $\Gamma^{\mu}_{\nu\rho}$}
The connection coefficients encoding the metric's affine structure:
\begin{equation}
\Gamma^{\mu}_{\nu\rho} = \frac{1}{2} g^{\mu\sigma} \left( \partial_\nu g_{\sigma\rho} + \partial_\rho g_{\sigma\nu} - \partial_\sigma g_{\nu\rho} \right)
\end{equation}

\subsection{Riemann Curvature Tensor $R^{\mu}_{\;\nu\rho\sigma}$}
Measures the intrinsic curvature of spacetime:
\begin{equation}
R^{\mu}_{\;\nu\rho\sigma} = \partial_\rho \Gamma^{\mu}_{\nu\sigma} - \partial_\sigma \Gamma^{\mu}_{\nu\rho} + \Gamma^{\mu}_{\rho\lambda} \Gamma^{\lambda}_{\nu\sigma} - \Gamma^{\mu}_{\sigma\lambda} \Gamma^{\lambda}_{\nu\rho}
\end{equation}

\subsection{Ricci Tensor $R_{\mu\nu}$ and Scalar $R$}
The contractions of the Riemann tensor:
\begin{equation}
R_{\mu\nu} = R^{\lambda}_{\;\mu\lambda\nu}, \quad R = g^{\mu\nu} R_{\mu\nu}
\end{equation}

\subsection{Einstein Tensor $G_{\mu\nu}$}
The curvature-matter coupling in GR:
\begin{equation}
G_{\mu\nu} = R_{\mu\nu} - \frac{1}{2} R g_{\mu\nu}
\end{equation}

\subsection{Geodesic Equations}
The equations governing free-fall motion:
\begin{equation}
\frac{d^2 x^\mu}{d\lambda^2} + \Gamma^{\mu}_{\nu\rho} \frac{dx^\nu}{d\lambda} \frac{dx^\rho}{d\lambda} = 0
\end{equation}

\subsection{Kretschmann Scalar $K$}
A quadratic curvature invariant useful for identifying singularities:
\begin{equation}
K = R^{\mu\nu\rho\sigma} R_{\mu\nu\rho\sigma}
\end{equation}

\section{Implementation Details}

\subsection{FLRW Metric}

The Friedmann-Lema\^{i}tre-Robertson-Walker metric for a flat universe is:
\begin{equation}
ds^2 = -c^2 dt^2 + a(t)^2 (dx^2 + dy^2 + dz^2)
\end{equation}

where $a(t)$ is the scale factor and $c$ is the speed of light. This metric describes a homogeneous and isotropic expanding universe.

\subsection{Schwarzschild Metric}

The Schwarzschild metric for a non-rotating black hole is:
\begin{equation}
ds^2 = -\left(1-\frac{2GM}{c^2 r}\right) c^2 dt^2 + \left(1-\frac{2GM}{c^2 r}\right)^{-1} dr^2 + r^2 d\theta^2 + r^2 \sin^2\theta d\phi^2
\end{equation}

where $G$ is Newton's gravitational constant and $M$ is the black hole mass.

\subsection{Geodesic Implementation}

The geodesic equations are derived from:
\begin{equation}
\ddot{x}^\mu + \Gamma^{\mu}_{\nu\rho} \dot{x}^\nu \dot{x}^\rho = 0
\end{equation}

where dots denote derivatives with respect to the affine parameter $\lambda$. The code computes symbolic expressions for each coordinate.

\subsection{Kretschmann Scalar Implementation}\label{sec:kretschmann}

The Kretschmann scalar is computed via:
\begin{equation}
K = R^{\alpha\beta\gamma\delta} R_{\alpha\beta\gamma\delta}
\end{equation}

where indices are raised using the inverse metric:
\begin{equation}
R^{\alpha\beta\gamma\delta} = g^{\alpha\mu} g^{\beta\nu} g^{\gamma\rho} g^{\delta\sigma} R_{\mu\nu\rho\sigma}
\end{equation}

For the Schwarzschild metric, this simplifies to the well-known result:
\begin{equation}
K = \frac{12 G^2 M^2}{c^4 r^6}
\end{equation}

For FLRW, the expression involves the scale factor and its derivatives:
\begin{equation}
K = \frac{12 \left(a^{2}{\left(t \right)} \left(\frac{d^{2}}{d t^{2}} a{\left(t \right)}\right)^{2} + \left(\frac{d}{d t} a{\left(t \right)}\right)^{4}\right)}{c^{4} a^{4}{\left(t \right)}}
\end{equation}

\section{Usage Examples}

\subsection{Basic Execution}

To compute tensors for the FLRW metric:
\begin{lstlisting}[language=bash]
python flrw_EinsteinKPy2Tex.py
\end{lstlisting}

For the Schwarzschild metric:
\begin{lstlisting}[language=bash]
python schw_EinsteinKPy2Tex.py
\end{lstlisting}

\subsection{Citation Management}

Display citation information:
\begin{lstlisting}[language=bash]
python einsteinpy_to_latex.py --cite
\end{lstlisting}

Generate BibTeX citation file:
\begin{lstlisting}[language=bash]
python flrw_EinsteinKPy2Tex.py --generate_bib
\end{lstlisting}

\subsection{\LaTeX\ Compilation}

Compile the generated \LaTeX\ document to PDF:
\begin{lstlisting}[language=bash]
pdflatex flrw_EinsteinKPy2Tex.tex
\end{lstlisting}

\section{Customizing for Your Own Metric}

The software is designed to be easily adaptable to any metric beyond the built-in FLRW and Schwarzschild cases. Users can compute tensors for custom spacetimes by modifying only two sections of the code.

\subsection{Modification Steps}

To compute tensors for a custom metric, the user modifies the following sections in the script:

\subsubsection*{1. Define coordinates and functions}

\begin{lstlisting}[language=Python]
# 1. Coordinates and functions
t, x, y, z = sp.symbols("t x y z")
coords = [t, x, y, z]
coord_names = ['t', 'x', 'y', 'z']  # for LaTeX labels

a = sp.Function("a")(t)  # Example: scale factor
c = sp.Symbol("c", positive=True, real=True)  # Speed of light
\end{lstlisting}

\subsubsection*{2. Define metric components}

\begin{lstlisting}[language=Python]
# 2. Metric (4D)
metric_arr = [
    [-c**2, 0, 0, 0],      # g_tt component
    [0, a**2, 0, 0],       # g_xx component
    [0, 0, a**2, 0],       # g_yy component
    [0, 0, 0, a**2],       # g_zz component
]
\end{lstlisting}

\subsection{Examples}

\subsubsection*{Example 1: Minkowski Metric (Flat Spacetime)}

The simplest spacetime, Minkowski metric, is obtained by setting $a(t) = 1$:

\begin{lstlisting}[language=Python]
# 1. Coordinates
t, x, y, z = sp.symbols("t x y z")
coords = [t, x, y, z]
c = sp.Symbol("c", positive=True, real=True)

# 2. Minkowski metric
metric_arr = [
    [-c**2, 0, 0, 0],
    [0, 1, 0, 0],
    [0, 0, 1, 0],
    [0, 0, 0, 1],
]
\end{lstlisting}

This yields zero Christoffel symbols, vanishing Riemann and Ricci tensors, $R = 0$, and $G_{\mu\nu} = 0$, as expected for flat spacetime.

\subsubsection*{Example 2: Custom Spherically Symmetric Metric}

For a general static, spherically symmetric metric of the form
\begin{equation}
ds^2 = -A(r) c^2 dt^2 + B(r) dr^2 + r^2 d\theta^2 + r^2 \sin^2\theta d\phi^2,
\end{equation}
the user defines:

\begin{lstlisting}[language=Python]
# 1. Coordinates
t, r, theta, phi = sp.symbols("t r theta phi")
coords = [t, r, theta, phi]
coord_names = ['t', 'r', '\\theta', '\\phi']

# 2. Define metric functions
A = sp.Function("A")(r)
B = sp.Function("B")(r)

# 3. Custom metric
metric_arr = [
    [-A(r), 0, 0, 0],
    [0, B(r), 0, 0],
    [0, 0, r**2, 0],
    [0, 0, 0, r**2 * sp.sin(theta)**2],
]
\end{lstlisting}

The software then automatically computes all tensors in terms of $A(r)$, $B(r)$, and their derivatives. This framework can accommodate any metric expressible in symbolic form, including Kerr (rotating black hole), Reissner-Nordstr\"om (charged black hole), or Bianchi (anisotropic cosmological) metrics.

\subsection{Important Considerations}

When defining a custom metric, users should carefully attend to several details to ensure correct tensor calculations and \LaTeX\ output.

First, \textbf{coordinate names} must be consistent between the coordinate list and the metric array indices. The software uses these names both for symbolic manipulation and for generating readable \LaTeX\ output. Mismatches can lead to silent errors or incorrectly labeled tensor components.

Second, \textbf{\LaTeX\ labels} in the \texttt{coord\_names} list require proper escaping for special characters. For example, the angular coordinate $\theta$ must be entered as \texttt{'$\backslash\backslash$theta'} so that the \LaTeX\ generator produces the correct math mode formatting. The same applies to $\phi$ (\texttt{'$\backslash\backslash$phi'}) and any other non-alphabetic symbols.

Third, the \textbf{metric signature} should be preserved according to the conventions of general relativity. The standard signature in most GR applications is $(-,+,+,+)$ (one timelike dimension and three spacelike dimensions). Altering this convention will propagate through all derived tensors, potentially leading to sign errors in physical quantities such as the Ricci scalar or Einstein tensor.

Fourth, \textbf{symbolic parameters} should be declared with appropriate mathematical assumptions using SymPy's symbol attributes. Specifying \texttt{positive=True} and \texttt{real=True} for physical quantities like mass $M$, gravitational constant $G$, speed of light $c$, and radial coordinate $r$ helps SymPy simplify expressions more aggressively, eliminating spurious absolute values and piecewise conditions. For example:
\begin{lstlisting}[language=Python]
M = sp.Symbol("M", positive=True, real=True)
r = sp.Symbol("r", positive=True, real=True)
\end{lstlisting}
This is particularly important for the Kretschmann scalar, where cancellations of terms like $(R_s - r)^2$ rely on SymPy's ability to assume $r > 0$.

Finally, users should verify that the \textbf{coordinate ordering} in the metric array matches the order of the coordinate list. The standard ordering in this software is $(t, x, y, z)$ for Cartesian coordinates and $(t, r, \theta, \phi)$ for spherical coordinates. The tensor computation routines assume this ordering throughout.

\subsection{Extending to Other Spacetimes}

The modular architecture of \href{https://github.com/lontelis/Einsteinpy_to_latex}{{\color{cyan}\texttt{EinsteinKPy2Tex}}} allows straightforward extension to any 4-dimensional metric beyond the two built-in examples. The user needs only to:

\begin{enumerate}
    \item \textbf{Define the coordinate system}: Choose appropriate coordinate symbols (e.g., $(t, x, y, z)$, $(t, r, \theta, \phi)$, or $(u, v, \theta, \phi)$ for null coordinates). The number of coordinates determines the spacetime dimension, which the software automatically propagates to all tensor calculations.

    \item \textbf{Declare metric functions and parameters}: Use SymPy's \texttt{Function} class for position-dependent metric components (e.g., $A(r)$, $B(r)$, $\Psi(t,r)$) or \texttt{Symbol} for constants (e.g., mass $M$, charge $Q$, cosmological constant $\Lambda$). The software will treat these symbolically, preserving them in the output expressions.

    \item \textbf{Construct the metric tensor array}: Create a $4 \times 4$ (or $D \times D$ for other dimensions) list of lists, where diagonal entries represent the metric components along each coordinate direction, and off-diagonal entries represent cross-terms. The array should be symmetric, as the metric tensor is symmetric by definition:
    \begin{lstlisting}[language=Python]
metric_arr = [
    [g_tt, g_tr, 0, 0],
    [g_tr, g_rr, 0, 0],
    [0, 0, g_thth, 0],
    [0, 0, 0, g_phph],
]
    \end{lstlisting}

    \item \textbf{Run the script} to obtain \LaTeX\ output. The software automatically computes all derived quantities without further user intervention.
\end{enumerate}

All subsequent tensor calculations—Christoffel symbols, Riemann curvature tensor, Ricci tensor and scalar, Einstein tensor, geodesic equations, and Kretschmann scalar—are handled automatically by the EinsteinPy backend \cite{Ribeiro2020} and the custom contraction routines described in Section~\ref{sec:kretschmann}. This automation eliminates manual algebra errors and ensures consistency across all derived quantities.

\section{Verification and Validation}

We verify our outputs against known analytical results from the general relativity literature. These tests confirm the correctness of both the EinsteinPy backend and our custom contraction routines.

\subsection{Schwarzschild Validation}

For the Schwarzschild metric, which describes the spacetime around a non-rotating, uncharged black hole, we perform the following checks:

\begin{itemize}
    \item \textbf{Ricci tensor}: The vacuum Einstein equations require $R_{\mu\nu} = 0$ outside the black hole. Our code returns zero for all components of the Ricci tensor, confirming the vacuum solution property.
    
    \item \textbf{Kretschmann scalar}: The fully contracted Riemann tensor product yields the well-known curvature invariant:
    \begin{equation}
    K = \frac{48 G^2 M^2}{c^4 r^6} = \frac{12 R_s^2}{r^6},
    \end{equation}
    where $R_s = 2GM/c^2$ is the Schwarzschild radius. This expression correctly diverges as $r \to 0$ (the physical singularity) while remaining finite at $r = R_s$ (the coordinate singularity at the event horizon).
    
    \item \textbf{Geodesic equations}: Our derived geodesic equations reproduce the standard orbital equations for massive and massless particles, including the perihelion precession of Mercury and light deflection by the Sun.
\end{itemize}

\subsection{FLRW Validation}

For the flat FLRW metric, which describes a homogeneous and isotropic expanding universe, we verify:

\begin{itemize}
    \item \textbf{Ricci scalar}: The scalar curvature simplifies to
    \begin{equation}
    R = \frac{6}{a^2}\left(\frac{\ddot{a}}{a} + \frac{\dot{a}^2}{c^2}\right),
    \end{equation}
    where $a(t)$ is the scale factor and dots denote derivatives with respect to cosmic time $t$. This matches standard cosmology textbooks.
    
    \item \textbf{Einstein tensor components}: The time-time component yields the Friedmann equation:
    \begin{equation}
    G_{tt} = 3\frac{\dot{a}^2}{c^2} = \frac{8\pi G}{c^4} \rho c^2,
    \end{equation}
    correctly relating the expansion rate to the energy density $\rho$.
    
    \item \textbf{Kretschmann scalar}: Our code produces
    \begin{equation}
    K = \frac{12\left(a^2 \ddot{a}^2 + \dot{a}^4\right)}{c^4 a^4},
    \end{equation}
    which reduces to zero for Minkowski spacetime ($a = \text{constant}$) as expected.
\end{itemize}

\subsection{Cross-Validation}

To ensure numerical stability, we also test the Schwarzschild metric in the weak-field limit ($r \gg R_s$), where it should approximate the Newtonian limit. The geodesic equations reduce to Newton's law of gravitation, and the Kretschmann scalar scales as $r^{-6}$, consistent with the inverse-square law of the gravitational field strength.

These validations confirm that \href{https://github.com/lontelis/Einsteinpy_to_latex}{{\color{cyan}\texttt{EinsteinKPy2Tex}}} produces correct results for both vacuum and cosmological spacetimes.

\section{Educational Applications}

The software is particularly well-suited for several educational and research contexts:

\begin{itemize}
    \item \textbf{GR Coursework}: Students can modify metric components and immediately see how the Christoffel symbols, Riemann tensor, and geodesic equations change. This interactive exploration deepens understanding of curvature and its relationship to the metric.

    \item \textbf{Research Quick-Checks}: Researchers can rapidly verify metric properties, such as whether a proposed metric satisfies Einstein's equations or whether it contains singularities (via the Kretschmann scalar). This accelerates the iterative process of metric construction.

    \item \textbf{Publication Output}: The direct \LaTeX\ export with non-zero components only allows authors to include tensor calculations directly in papers without manual typesetting. This reduces transcription errors and saves significant time.

    \item \textbf{Self-Study}: Graduate students and self-learners can visualize which tensor components are non-zero and how they depend on coordinates, building intuition for curved spacetimes.

    \item \textbf{Code Reuse}: The modular design allows users to extract specific calculation routines for integration into larger projects, such as numerical relativity codes or cosmological parameter estimation pipelines.
\end{itemize}

\section{Comparison with Existing Tools}

In \autoref{tab:comparison}, we present a comparison with existing tools. The Kretschmann scalar ($K$) is a key feature highlighted by the ``K'' in the name. GRTensor (a Maple package) and xAct (a Mathematica suite) are powerful but require commercial software licenses. Base EinsteinPy provides the foundational tensors but lacks geodesic equations, Kretschmann scalar, and \LaTeX\ export. \href{https://github.com/lontelis/Einsteinpy_to_latex}{{\color{cyan}\texttt{EinsteinKPy2Tex}}} uniquely combines Python-based symbolic computation with direct \LaTeX\ export of non-zero components, geodesic equations, and the Kretschmann scalar, all under an open-source license.

\begin{table}[ht!]
%\centering
%\hspace{-1cm}
\begin{tabular}{lcccc}
\toprule
\textbf{Feature} & \textbf{EinsteinKPy2Tex} & \textbf{GRTensor} & \textbf{xAct} & \textbf{EinsteinPy} \\
\midrule
Open source & \checkmark & \texttimes & \checkmark & \checkmark \\
Python-based & \checkmark & \texttimes & \texttimes & \checkmark \\
LaTeX export & \checkmark & Limited & \checkmark & \texttimes \\
Geodesic equations & \checkmark & \checkmark & \checkmark & \texttimes \\
\textbf{Kretschmann scalar} & \textbf{\checkmark} & \checkmark & \checkmark & \texttimes \\
Built-in citations & \checkmark & \texttimes & \texttimes & \texttimes \\
Zero-cost & \checkmark & \texttimes & \checkmark & \checkmark \\
\midrule
\textbf{Platform} & Python & Maple & Mathematica & Python \\
\midrule
\end{tabular}
\caption{Comparison of \href{https://github.com/lontelis/Einsteinpy_to_latex}{{\color{cyan}\texttt{EinsteinKPy2Tex}}} with existing GR tensor calculation tools. }
\label{tab:comparison}
\end{table}

\section{Numerical Simulation of Kretschmann Scalar}

To validate our symbolic computations and visualize the physical implications, 
we perform numerical simulations of the Kretschmann scalar for both metrics 
across a range of parameters. All simulations use natural units with 
$G = c = 1$ for simplicity.

\subsection{Schwarzschild Black Holes}

The Kretschmann scalar for a Schwarzschild black hole is:
\begin{equation}
K(r) = \frac{12 R_s^2}{r^6}, \quad R_s = 2M
\end{equation}

Figure~\ref{fig:schwarzschild} shows $K(r)$ for different black hole masses 
$M = 0.5, 1.0, 2.0, 5.0$ (in solar mass units). The left panel displays 
the linear-scale behavior: $K(r)$ decreases monotonically with increasing 
radius and is larger for more massive black holes. The right panel uses a 
log-log scale, clearly demonstrating the $K \propto r^{-6}$ power-law 
behavior (reference line shown). The event horizon at $r = 2M$ is marked 
for each mass; note that $K(r)$ remains finite at the horizon, confirming 
it is a coordinate singularity rather than a physical one.

The simulations confirm the physical singularity at $r=0$ (where $K \to 
\infty$) and the $r^{-6}$ falloff at large distances, consistent with the 
inverse-square law of the gravitational field strength.

\begin{figure}[H]
    \centering
    \includegraphics[width=0.9\textwidth]{schwarzschild_kretschmann.pdf}
    \caption{Kretschmann scalar for Schwarzschild metric with different 
             black hole masses. Left: Linear scale showing amplitude dependence 
             on mass. Right: Log-log scale confirming $K \propto r^{-6}$ 
             power-law behavior (dashed reference line). The event horizon 
             $r_s = 2M$ is indicated for each mass.}
    \label{fig:schwarzschild}
\end{figure}

\subsection{FLRW Cosmological Models}

For the FLRW metric, the Kretschmann scalar is:
\begin{equation}
K(t) = 12\left(\frac{\ddot{a}^2}{a^2} + \frac{\dot{a}^4}{a^4}\right)
\end{equation}
where $a(t)$ is the scale factor. We simulate three distinct cosmological eras:

\begin{itemize}
    \item \textbf{Radiation-dominated}: $a(t) \propto t^{1/2}$, giving 
          $K(t) \propto t^{-4}$
    \item \textbf{Matter-dominated}: $a(t) \propto t^{2/3}$, giving 
          $K(t) \propto t^{-4}$
    \item \textbf{$\Lambda$-dominated (de Sitter)}: $a(t) \propto e^{Ht}$, 
          giving $K(t) \propto \text{constant}$.
\end{itemize}

Figure~\ref{fig:flrw} presents the Kretschmann scalar evolution for these 
three scenarios. The left panel shows the early universe behavior (linear 
scale), where the radiation-dominated era exhibits the highest curvature. 
The right panel uses a log-log scale, confirming the predicted power-law 
scalings: $K \propto t^{-4}$ for radiation domination and $K \propto t^{-4}$ 
for matter domination (reference lines shown). For $\Lambda$-domination, 
$K$ decays exponentially, reflecting the de Sitter horizon.

Notably, all three models exhibit a big bang singularity at $t=0$ where 
$K \to \infty$, consistent with the initial singularity of standard 
cosmology.

\begin{figure}[H]
    \centering
    \includegraphics[width=0.9\textwidth]{flrw_kretschmann.pdf}
    \caption{Kretschmann scalar evolution for FLRW cosmologies. Left: Early 
             universe behavior showing radiation domination (red) producing 
             the highest curvature. Right: Log-log scale confirming 
             $K \propto t^{-4}$ (radiation) and $K \propto t^{-4}$ (matter) 
             power laws. The Kretschmann scalar evolution for the $\Lambda$-dominated case (green) is constant, as a function of cosmic time.}
    \label{fig:flrw}
\end{figure}

\subsection{Singularity Comparison}

Figure~\ref{fig:singularities} directly compares the singularity structure 
of both spacetimes. For Schwarzschild (left), the Kretschmann scalar diverges 
as $r \to 0$ (physical singularity) but remains finite at the event horizon 
$r = 2M$. For FLRW (right), the divergence occurs as $t \to 0$ (big bang 
singularity), with the radiation-dominated universe diverging steeply 
($t^{-4}$) and similarly for the the matter-dominated ($t^{-4}$).

\begin{figure}[H]
    \centering
    \includegraphics[width=0.9\textwidth]{kretschmann_singularities.pdf}
    \caption{Comparison of singularities. Left: Schwarzschild black hole 
             showing singularity at $r=0$ and finite $K$ at the event horizon. 
             Right: FLRW universe showing big bang singularity at $t=0$ for 
             different cosmological eras.}
    \label{fig:singularities}
\end{figure}

\subsection{3D Parameter Space Exploration}

Figure~\ref{fig:schwarzschild_3d} extends the Schwarzschild analysis to a 
3D parameter space, showing $K(r,M)$ as a function of both radius and mass. 
The surface plot demonstrates that increasing mass increases the Kretschmann 
scalar amplitude for fixed $r$, while decreasing $r$ increases $K$ with a 
steeper dependence ($r^{-6}$).

\begin{figure}[H]
    \centering
    \includegraphics[width=0.8\textwidth]{schwarzschild_3d.pdf}
    \caption{Three-dimensional visualization of the Kretschmann scalar 
             $K(r,M)$ for Schwarzschild black holes. The surface shows the 
             combined dependence on radial coordinate and black hole mass.}
    \label{fig:schwarzschild_3d}
\end{figure}

\subsection{Simulation Summary}

These numerical simulations confirm that EinsteinKPy2Tex correctly computes 
the Kretschmann scalar for both metrics. The results reproduce:
\begin{itemize}
    \item The finite curvature at the Schwarzschild event horizon
    \item The physical singularity at $r=0$ for Schwarzschild
    \item The big bang singularity at $t=0$ for FLRW
    \item The curvature from the correct power-law scalings for Kretschmann scalar are the following: \\ $r^{-6}$ (Schwarzschild) and $t^{-4}/t^{-4}/\text{constant}$ (FLRW)
    \item The curvature from Kretschmann scalar remains constant for $\Lambda$-dominated universes.
\end{itemize}

\newpage %%% %%% %%% %%% %%% %%% %%% %%%

\section{Physical Interpretation of the txyzS Metric and Its Kretschmann Scalar}

\subsection{Motivation from Advanced Manifold-Metric Pairs Theory}

The txyzS metric is motivated by the AMMP theory framework introduced in Ntelis (2025) \cite{Ntelis2025_AMMP}. In this framework, physical quantities traditionally treated as parameters (such as entropy, temperature, or mass) can be promoted to genuine spacetime coordinates, provided they satisfy appropriate metric compatibility conditions. The entropy coordinate $S$ is a natural candidate for such an extension due to its fundamental role in thermodynamics and information theory, as well as its connection to spacetime geometry via the Bekenstein-Hawking formula $S = k_B A/(4\ell_P^2)$.

For the metric with entropy as an extra dimension,
\begin{equation}
ds^2 = -c^2 dt^2 + a(t)^2 (dx^2+dy^2+dz^2) + f_S(t,S)^2 dS^2,
\end{equation}
the Kretschmann scalar takes the relatively simple form:

\begin{align}
K = \frac{4}{c^4}\left(\frac{\ddot{f}_S}{f_S}\right)^2 
+ \frac{12}{c^4}\left(\frac{\ddot{a}}{a}\right)^2 
+ \frac{12}{c^4}\left(\frac{\dot{a}}{a}\right)^2\left(\frac{\dot{f}_S}{f_S}\right)^2 
+ \frac{12}{c^4}\left(\frac{\dot{a}}{a}\right)^4.
\end{align}
where dots denote derivatives with respect to cosmic time $t$, and 
$\dot{f}_S = \partial f_S/\partial t$ at fixed entropy $S$.

\subsection{Choice of the Entropy Scale Factor $f_S(t,S)$}

The functional form of $f_S(t,S)$ must reflect the physical behavior of entropy in gravitational systems. We choose the exponential form
\begin{equation}
f_S(t,S) = e^{\beta t S},
\end{equation}
where $\beta$ is a constant parameter with dimensions of inverse action (in natural units, $\beta$ is dimensionless). This choice is motivated by several physical considerations grounded in the AMMP framework:

\begin{enumerate}
    \item \textbf{Second Law of Thermodynamics}: The second law requires 
          $\dot{S} \geq 0$ for isolated systems. The exponential form 
          satisfies $\partial f_S/\partial t = \beta S f_S > 0$ for 
          $t > 0$ and $S > 0$, ensuring entropy never decreases.

\item \textbf{Cosmological Entropy Production}: In the expanding universe, 
          the entropy within the cosmological horizon scales with the 
          horizon area. The exponential form $f_S(t,S) = e^{\beta t S}$ 
          captures this growth while maintaining a separable dependence 
          on $t$ and $S$. For fixed entropy coordinate $S$, the scale 
          factor grows exponentially with time, analogous to the de Sitter 
          expansion of spatial dimensions. The parameter $\beta$ controls 
          the expansion rate of the entropy dimension relative to the 
          spatial dimensions.

\item \textbf{Separability and Dimensional Analysis}: The product $\beta t S$ 
          is dimensionless in natural units, ensuring $f_S$ is well-defined. 
          The exponential form is the unique separable solution (up to 
          rescaling) that yields constant fractional growth rates:
          $\dot{f}_S/f_S = \beta S$ and $f_S'/f_S = \beta t$, where prime 
          denotes derivative with respect to $S$. This self-similarity 
          simplifies the Kretschmann scalar significantly.

    \item \textbf{Computational Tractability}: The exponential form yields 
          simple derivatives:
          \begin{align}
          \dot{f}_S &= \beta S f_S, \\
          \ddot{f}_S &= (\beta S)^2 f_S,
          \end{align}
          which simplify the Kretschmann scalar to:
          \begin{align}
          K_{\text{txyzS}} = \frac{4}{c^4}(\beta S)^4 
          + \frac{12}{c^4}\left(\frac{\ddot{a}}{a}\right)^2 
          + \frac{12}{c^4}\left(\frac{\dot{a}}{a}\right)^2 (\beta S)^2 
          + \frac{12}{c^4}\left(\frac{\dot{a}}{a}\right)^4.
          \end{align}
          This decomposition separates the entropy contribution ($\beta S$) 
          from the cosmological dynamics ($a(t)$), allowing clear physical 
          interpretation.

    \item \textbf{AMMP Consistency}: The AMMP framework requires that the 
          metric functions satisfy certain compatibility conditions. The 
          exponential form ensures that the extra dimension remains 
          non-degenerate ($f_S > 0$) and that the metric signature is 
          preserved for all $t$ and $S$.

    \item \textbf{Thermodynamic Limit}: For small $\beta t S$, the exponential 
          reduces to $f_S \approx 1 + \beta t S$, recovering a linear entropy 
          dependence. The linear regime approximates early-time or low-entropy 
          systems, while the exponential regime dominates at late times or 
          high entropy.
\end{enumerate}

For the simulations presented in this work, we set $\beta = 0.5$ in natural 
units. This value is chosen to produce deviations from FLRW of order unity 
for $S \sim 1$, making the entropy dimension effects potentially detectable. 
Future work will constrain $\beta$ using observational data from the cosmic 
microwave background and large-scale structure.

\subsection{Mathematical expression comparison between $K_{\text{txyzS}}$ and $K_{\text{FLRW}}$}\label{sec:kretschmann_comparison}
We found that Kretschmann scalar for FLRW metric is 
          \begin{align}
          K_{\text{FLRW}} = \frac{12}{c^4}\left(\frac{\ddot{a}}{a}\right)^2 
          + \frac{12}{c^4}\left(\frac{\dot{a}}{a}\right)^4.
          \end{align}

We found that Kretschmann scalar for txyzS metric is 
          \begin{align}
          K_{\text{txyzS}} = \frac{4}{c^4}(\beta S)^4 
          + \frac{12}{c^4}\left(\frac{\ddot{a}}{a}\right)^2 
          + \frac{12}{c^4}\left(\frac{\dot{a}}{a}\right)^2 (\beta S)^2 
          + \frac{12}{c^4}\left(\frac{\dot{a}}{a}\right)^4.
          \end{align}
          
So the ratio is           

          \begin{align}
          \frac{ K_{\text{txyzS}} }{ K_{\text{FLRW}} } 
          = 
          1 + 
          \frac{
          \frac{4}{c^4}(\beta S)^4 
          + 
          \frac{12}{c^4}\left(\frac{\dot{a}}{a}\right)^2 (\beta S)^2 
          }{
          \frac{12}{c^4}\left(\frac{\ddot{a}}{a}\right)^2 
          + \frac{12}{c^4}\left(\frac{\dot{a}}{a}\right)^4
          }
          \end{align}

\subsection{Physical Interpretation of Kretschmann Terms}

The decomposition of $K$ admits a clear physical interpretation:

\begin{itemize}
    \item \textbf{$\left(\ddot{a}/a\right)^2$ term}: Measures the acceleration 
          of the cosmic expansion. In standard cosmology, this relates to 
          the deceleration parameter $q = -\ddot{a}a/\dot{a}^2$.

    \item \textbf{$\left(\dot{a}/a\right)^4$ term}: The square of the Hubble 
          parameter $H = \dot{a}/a$, representing the expansion rate. This term
          dominates in the late universe (dark energy era).

    \item \textbf{$\left(\ddot{f}_S/f_S\right)^2 = (\beta S)^4$ term}: The 
          "entropy acceleration" — how the entropy scale factor changes its 
          rate of change. In black hole thermodynamics, this relates to the 
          second law $dS/dt \geq 0$, with $\ddot{f}_S$ representing the 
          entropy production rate. For the exponential ansatz, this term is 
          constant in time, implying steady entropy production.

    \item \textbf{$\left(\dot{f}_S/f_S\right)^4 = (\beta S)^4$ term}: The 
          "entropy Hubble parameter" $H_S = \dot{f}_S/f_S = \beta S$, 
          representing the fractional rate of entropy increase. For 
          Schwarzschild black holes, $H_S$ would be related to the Hawking 
          evaporation rate.

    \item \textbf{Coupling term $\left(\dot{a}/a\right)^2 \left(\dot{f}_S/f_S\right)^2 = (\beta S)^2 H^2$}: 
          Represents the interaction between cosmic expansion and entropy 
          production. This term is unique to the 5D metric and has no 
          counterpart in 4D GR, potentially describing how the universe's 
          expansion drives entropy increase (or vice versa). For $\beta S = 1$ 
          and $H = 1$, this term contributes equally to $K$ as the $H^4$ term.
\end{itemize}

This decomposition makes the txyzS metric computationally tractable 
and physically interpretable. A full 7D simulation with time, $t$, spatial space, $(x,y,z)$, mass $M$, temperature 
$T$, and entropy $S$ would produce a Kretschmann scalar with over 10 terms, 
complicating analysis unnecessarily, at this stage. The entropy-only case, i.e. the txyzS metric case, retains the 
essential thermodynamic dimension while remaining analytically and 
numerically manageable.

\section{Numerical Simulation of the Entropy Dimension}

\subsection{Simulation Setup}

We simulate the Kretschmann scalar for both the standard FLRW metric and our 
txyzS metric with entropy as an extra dimension. For the FLRW case,
\begin{equation}
K_{\text{FLRW}} = \frac{12}{c^4}\left[\left(\frac{\ddot{a}}{a}\right)^2 + \left(\frac{\dot{a}}{a}\right)^4\right].
\end{equation}
For the txyzS metric with the exponential entropy ansatz $f_S(t,S) = e^{\beta t S}$,
the Kretschmann scalar simplifies to:
\begin{equation}
K_{\text{txyzS}} = \frac{12}{c^4}\left(\frac{\ddot{a}}{a}\right)^2 
+ \frac{12}{c^4}\left(\frac{\dot{a}}{a}\right)^4
+ \frac{4}{c^4}(\beta S)^4
+ \frac{12}{c^4}\left(\frac{\dot{a}}{a}\right)^2 (\beta S)^2.
\end{equation}

We assume natural units $c = 1$, set $\beta = 0.5$, and consider three 
cosmological eras:
\begin{itemize}
    \item Matter-dominated: $a(t) \propto t^{2/3}$
    \item Radiation-dominated: $a(t) \propto t^{1/2}$
    \item $\Lambda$-dominated (de Sitter): $a(t) \propto e^{Ht}$ with $H=1$
\end{itemize}
The entropy coordinate $S$ is varied from $0$ to $2$ to explore the 
parameter space.

\subsection{Results}

Figure~\ref{fig:flrw} shows the FLRW Kretschmann scalar for the three 
cosmological eras. The power-law scalings $K \propto t^{-4}$ (radiation) 
and $K \propto t^{-4}$ (matter) are clearly visible, as expected from 
analytic theory. The $\Lambda$-dominated case shows exponential decay 
of curvature at late times.

\begin{figure}[H]
    \centering
    \includegraphics[width=0.8\textwidth]{fig1_flrw_kretschmann.pdf}
    \caption{FLRW Kretschmann scalar for different cosmological eras. 
             The radiation-dominated era exhibits $K \propto t^{-4}$, 
             matter-dominated shows $K \propto t^{-4}$, while the 
             $\Lambda$-dominated universe approaches a constant curvature 
             at all cosmic times.}
    \label{fig:flrw}
\end{figure}

\autoref{fig:txyzS_vs_flrw} compares the txyzS metric (with entropy 
dimension) to standard FLRW in a matter-dominated universe. For increasing 
entropy coordinate $S$, the Kretschmann scalar increases, indicating that 
the entropy dimension amplifies spacetime curvature. The amplification is 
most pronounced at early times ($t \ll 1$), where the $(\beta S)^4$ term 
dominates over the $H^4$ term.

\begin{figure}[H]
    \centering
    \includegraphics[width=0.8\textwidth]{fig2_txyzS_vs_flrw_matter.pdf}
    \caption{Comparison of Kretschmann scalar: FLRW vs txyzS (matter-dominated). 
             The entropy dimension introduces an additional constant term 
             $4(\beta S)^4$ and a coupling term $12H^2(\beta S)^2$, both 
             increasing $K$ relative to FLRW.}
    \label{fig:txyzS_vs_flrw}
\end{figure}

\iffalse
Figure~\ref{fig:ratio} shows the ratio $K_{\text{txyzS}}/K_{\text{FLRW}}$ 
for different $S$ values. At late times ($t \gg 1$), the ratio approaches 
$1 + (\beta S)^4 / H^4 + 3(\beta S)^2 / H^2$, which for $H \to 0$ in an 
expanding universe tends to infinity. At early times ($t \ll 1$), the 
ratio is dominated by the constant entropy term, giving $R \approx 
1 + 4(\beta S)^4 / (12H^4) \approx 1 + (\beta S)^4/(3H^4)$. For 
$S = 2$, $\beta = 0.5$, and $H \sim 1$, this yields $R \approx 1.33$.

\begin{figure}[H]
    \centering
    \includegraphics[width=0.8\textwidth]{fig4_ratio_comparison.pdf}
    \caption{Ratio $K_{\text{txyzS}}/K_{\text{FLRW}}$ for different entropy 
             coordinates $S$ in a matter-dominated universe. The ratio 
             exceeds 1 for all $S > 0$, with larger $S$ producing larger 
             deviations from FLRW. At early times ($t \ll 1$), the ratio 
             approaches a constant determined by $(\beta S)^4 / H^4$.}
    \label{fig:ratio}
\end{figure}
\fi


%%%%%%%%%%%%%%%%%%%%%%%%%%

\newpage

\autoref{fig:ratio} shows the ratio $K_{\text{txyzS}}/K_{\text{FLRW}}$ 
for different $S$ values. From the explicit expressions derived in 
Section~\ref{sec:kretschmann_comparison}, the ratio is given by:

\begin{align}
\frac{K_{\text{txyzS}}}{K_{\text{FLRW}}} = 
1 + \frac{
\frac{4}{c^4}(\beta S)^4 
+ \frac{12}{c^4}\left(\frac{\dot{a}}{a}\right)^2 (\beta S)^2
}{
\frac{12}{c^4}\left(\frac{\ddot{a}}{a}\right)^2 
+ \frac{12}{c^4}\left(\frac{\dot{a}}{a}\right)^4
}.
\end{align}

This expression reveals several key features:

\begin{itemize}
    \item \textbf{Constant numerator contribution}: The term $\frac{4}{c^4}(\beta S)^4$ 
          is independent of cosmic time $t$ and depends only on the entropy 
          coordinate $S$ and the AMMP parameter $\beta$. This ensures that 
          the ratio never reduces to exactly 1 for any $S > 0$, even in the 
          limit of vanishing expansion rates.

    \item \textbf{Coupling term}: The term $\frac{12}{c^4}\left(\frac{\dot{a}}{a}\right)^2 (\beta S)^2$ 
          couples the Hubble parameter $H = \dot{a}/a$ with the entropy 
          dimension. At early times ($t \ll 1$), $H$ is large, making this 
          term dominate the numerator alongside the constant term.

    \item \textbf{Denominator behavior}: The denominator contains both 
          $(\ddot{a}/a)^2$ and $H^4$. For matter-dominated universes 
          ($a \propto t^{2/3}$), both terms scale as $t^{-4}$, while for 
          radiation-dominated universes ($a \propto t^{1/2}$), they also scale 
          as $t^{-4}$.

    \item \textbf{Late-time behavior}: As $t \to \infty$ in a matter-dominated 
          or $\Lambda$-dominated universe, $H \to 0$ and $\ddot{a}/a \to 0$, 
          so the denominator approaches zero. Consequently, the ratio 
          diverges ($\to \infty$) at late times, indicating that the entropy 
          dimension dominates the curvature invariant in the distant future.

    \item \textbf{Early-time behavior}: As $t \to 0$ (early universe), 
          $H \to \infty$ and $\ddot{a}/a \to \infty$, with both scaling as 
          power laws. For matter domination ($H \sim 2/(3t)$, $\ddot{a}/a \sim -2/(9t^2)$), 
          the denominator scales as $t^{-4}$. The numerator contains a 
          constant term $(\beta S)^4$ and a term scaling as $t^{-2}$ (since 
          $H^2 \propto t^{-2}$). Therefore, at sufficiently early times, 
          the $t^{-2}$ term dominates, and the ratio approaches a constant 
          determined by the relative coefficients.
\end{itemize}

For the specific case of matter-dominated universe ($a \propto t^{2/3}$), 
we have $H = 2/(3t)$ and $\ddot{a}/a = -2/(9t^2)$. Substituting into the 
ratio expression (with $c=1$ in natural units):

\begin{align}
\frac{K_{\text{txyzS}}}{K_{\text{FLRW}}} = 
1 + \frac{
4(\beta S)^4 + 12\left(\frac{2}{3t}\right)^2 (\beta S)^2
}{
12\left(-\frac{2}{9t^2}\right)^2 + 12\left(\frac{2}{3t}\right)^4
}
= 1 + \frac{
4(\beta S)^4 + \frac{16}{3t^2} (\beta S)^2
}{
\frac{16}{27t^4} + \frac{64}{27t^4}
}
= 1 + \frac{
4(\beta S)^4 + \frac{16}{3t^2} (\beta S)^2
}{
\frac{80}{27t^4}
}.
\end{align}

Simplifying:

\begin{align}
\frac{K_{\text{txyzS}}}{K_{\text{FLRW}}} = 
1 + \frac{27t^4}{80} \left[4(\beta S)^4 + \frac{16}{3t^2} (\beta S)^2\right]
= 1 + \frac{27}{20} t^4 (\beta S)^4 + \frac{9}{5} t^2 (\beta S)^2.
\end{align}

Thus:

\begin{align}
\frac{K_{\text{txyzS}}}{K_{\text{FLRW}}} 
= 1 + \frac{27}{20} t^4 (\beta S)^4 + \frac{9}{5} t^2 (\beta S)^2.
\end{align}

For $\beta = 0.5$, $S = 2$ (so $\beta S = 1$), and $t = 1$ (present time in 
natural units), this yields:

\begin{align}
\frac{K_{\text{txyzS}}}{K_{\text{FLRW}}} 
%= 1 + \frac{27}{20}  (1)^4 + \frac{9}{5} (1)^2 = \frac{83}{20} = 4.15 .
= \frac{83}{20} = 4.15 .
\end{align}

At earlier times ($t = 0.1$), the ratio is:

\begin{align}
\frac{K_{\text{txyzS}}}{K_{\text{FLRW}}} 
%= 1 + \frac{27}{20}  (0.0001) + \frac{9}{5} (0.01) = \frac{203627}{200000} = 1.018135.
= \frac{203627}{200000} = 1.018135.
\end{align}

This demonstrates that the entropy dimension produces the largest 
deviations from FLRW at late times (present and future) rather than 
early times, contrary to naive expectations. The constant term 
$4(\beta S)^4$ in the numerator becomes increasingly significant 
relative to the denominator as the universe expands and curvature 
decreases.

\begin{figure}[H]
    \centering
    \includegraphics[width=0.8\textwidth]{fig4_ratio_comparison.pdf}
    \caption{Ratio $K_{\text{txyzS}}/K_{\text{FLRW}}$ for different entropy 
             coordinates $S$ in a matter-dominated universe with $\beta = 0.5$. 
             The ratio exceeds 1 for all $S > 0$ and diverges at late times 
             ($t \to \infty$) as the denominator approaches zero. The analytical 
             expression derived in the text matches the numerical simulation 
             exactly.}
    \label{fig:ratio}
\end{figure}

%%%%%%%%%%%%%%%%%%%%%%%%%%


\subsection{Discussion of Exponential Entropy Ansatz}

The choice $f_S(t,S) = e^{\beta t S}$ is not arbitrary; it emerges naturally 
from the AMMP framework \cite{Ntelis2025_AMMP} when requiring that the entropy 
dimension be compatible with the second law of thermodynamics and the 
evolution equations for the scale factor. The exponential form has several 
desirable properties:

\begin{enumerate}
    \item \textbf{Separability}: The dependence on $t$ and $S$ separates, 
          allowing analytic simplification of the Kretschmann scalar.

    \item \textbf{Dimensionless argument}: The product $\beta t S$ is 
          dimensionless in natural units, ensuring $f_S$ is well-defined 
          regardless of the units of $t$ and $S$.

    \item \textbf{Monotonicity}: For $t > 0$ and $S > 0$, $f_S$ increases 
          with both arguments, consistent with the second law.

    \item \textbf{Constant entropy Hubble parameter}: $H_S = \dot{f}_S/f_S = \beta S$ 
          is constant in time, implying steady entropy production. This is 
          consistent with the observed approximately constant entropy 
          production rate in the universe's expansion.

    \item \textbf{Perturbative limit}: For small $\beta t S$, the exponential 
          reduces to the linear regime $f_S \approx 1 + \beta t S$, allowing 
          connection to simpler entropy models.
\end{enumerate}

The simulation results demonstrate that the entropy dimension produces 
observable deviations from standard FLRW cosmology, particularly at late 
times and for larger $S$. Future work will constrain $\beta$ using 
observational data from the cosmic microwave background, baryon acoustic 
oscillations, and supernova surveys. The AMMP framework provides a 
theoretical foundation for extending these calculations to include 
temperature and mass dimensions, which will be explored in subsequent 
publications.

\subsection{Discussion}

The entropy dimension introduces three new effects:

\begin{enumerate}
    \item \textbf{Direct entropy contribution}: The term $(\ddot{f}_S/f_S)^2$ 
          represents the acceleration of entropy growth. For exponential 
          entropy growth, this term is constant in time.

    \item \textbf{Coupling term}: The cross-term $(\dot{a}/a)^2(\dot{f}_S/f_S)^2$ 
          couples cosmic expansion with entropy production. This term is 
          unique to theories with an entropy dimension.

    \item \textbf{Amplification of curvature}: For $S > 0$, the entropy 
          dimension increases the overall Kretschmann scalar, suggesting 
          that entropy gradients contribute to spacetime curvature.
\end{enumerate}

The strongest deviation from FLRW occurs in the late universe ($t > 0.1$), 
where the ratio $K_{\text{txyzS}}/K_{\text{FLRW}}$ can exceed 2 for 
$S = 2$. This suggests that entropy dimension effects would be most 
detectable in the early universe, potentially leaving signatures in the 
cosmic microwave background.

\subsection{Conclusion of Simulation}

The txyzS metric is computationally tractable and produces 
physically interpretable results. The simulation confirms that:
\begin{itemize}
    \item The entropy dimension amplifies curvature at early times
    \item The effect scales with the entropy coordinate $S$
    \item The coupling between expansion and entropy is unique to this model
\end{itemize}
Future work will explore observational constraints from CMB data and 
black hole thermodynamics.

\section{Discussion for future work}

Several enhancements are planned for future releases of \href{https://github.com/lontelis/Einsteinpy_to_latex}{{\color{cyan}\texttt{EinsteinKPy2Tex}}}:

\begin{itemize}
    \item \textbf{Kerr metric}: Support for rotating black holes, including the Kerr metric and the Kerr-Newman generalization with charge. This requires handling off-diagonal metric components ($g_{t\phi} \neq 0$) and more complex geodesic equations.

    \item \textbf{FLRW with curvature}: Extension to non-flat FLRW metrics with spatial curvature $k = \pm 1$ (closed and open universes). This involves additional metric components and modifications to the Kretschmann scalar.

    \item \textbf{Command-line metric definition}: A user-friendly interface allowing custom metric definitions without modifying the source code, perhaps via a configuration file or interactive prompts.

    \item \textbf{Jupyter notebook integration}: Interactive notebooks with live visualization of tensor components, geodesic trajectories, and curvature invariants as functions of coordinates.

    \item \textbf{Plotting capabilities}: Automatic generation of plots for the Kretschmann scalar as a function of radius, showing the singularity structure of black hole spacetimes.

    \item \textbf{Web-based interface}: A browser-based interface for non-Python users, enabling metric definition and tensor computation without local installation.
\end{itemize}

\subsection{Long-term prospects}

While \href{https://github.com/lontelis/Einsteinpy_to_latex}{{\color{cyan}\texttt{EinsteinKPy2Tex}}} currently focuses on local differential geometric quantities, future extensions could include:

\begin{itemize}
    \item \textbf{Topological invariants}: Computation of the Euler characteristic $\chi(M)$ and Chern classes for compact spacetimes, using the Gauss-Bonnet theorem and characteristic class formulas. This would require integration over the manifold, not just local tensor evaluation.

    \item \textbf{Gauss-Bonnet term}: Integration of the Gauss-Bonnet density $\sqrt{|g|}(R^2 - 4R_{\mu\nu}R^{\mu\nu} + R_{\mu\nu\rho\sigma}R^{\mu\nu\rho\sigma})$ for studying modified gravity theories.    

    \item \textbf{Pontryagin classes}: Calculation of Pontryagin classes for gravitational instantons and other topologically non-trivial spacetimes in Euclidean quantum gravity.

    \item \textbf{AdS/CFT applications}: Computation of Weyl tensor invariants and conformal anomalies relevant to the AdS/CFT correspondence.

    \item \textbf{Functors of actions theories term}: Integration of the functors of actions theories \cite{Ntelis2023} for studying extended gravity and cosmology theories.    

    \item \textbf{Advanced manifold-metric pairs}: Extensive studies of extra dimensions from advanced manifold-metric pairs theories \cite{Ntelis2025_AMMP} for studying extended gravity and cosmology theories.    

\end{itemize}

However, these topological extensions require global manifold information (compactness, boundaries, orientation) and integration routines that are beyond the current scope of \href{https://github.com/lontelis/Einsteinpy_to_latex}{{\color{cyan}\texttt{EinsteinKPy2Tex}}}. They remain an interesting direction for future research and collaboration.

\section{Conclusions}

\href{https://github.com/lontelis/Einsteinpy_to_latex}{{\color{cyan}\texttt{EinsteinKPy2Tex}}} provides a reliable, user-friendly bridge between symbolic GR calculations and publication-ready \LaTeX. By automating tensor computations for the most important metrics in cosmology and black hole physics, it saves researchers and educators significant time while reducing errors. The built-in citation system ensures proper academic attribution, and the MIT license encourages wide adoption and contribution.

Beyond its utility as a computational tool, \href{https://github.com/lontelis/Einsteinpy_to_latex}{{\color{cyan}\texttt{EinsteinKPy2Tex}}} 
addresses a 
growing research trend that seeks to unify black hole and cosmological 
physics through common curvature invariants. While the Kretschmann scalar 
has been known for decades, our software provides the first unified 
platform to compute, verify, and visualize these invariants side-by-side 
for the two most important metrics in general relativity. This capability 
enables researchers to explore proposed connections---such as the 
geometric bridge proposed by Sussman \& Hernandez (2019) \cite{Sussman2019}---and offers 
students an interactive window into the singularities that define our 
universe. As we extend EinsteinKPy2Tex to include rotating (Kerr) and 
charged (Reissner-Nordstr\"om) black holes, as well as curved FLRW 
cosmologies, we anticipate its role in both pedagogy and research will 
continue to grow. The software is freely available, fully reproducible, 
and welcomes community contributions.

In this paper, we have presented the first concrete application of the 
Advanced Manifold-Metric Pairs (AMMP) framework developed by Ntelis (2025) \cite{Ntelis2025_AMMP}. 
By promoting entropy from a thermodynamic parameter to a genuine spacetime 
coordinate, we constructed a novel 5-dimensional metric—the txyzS metric—that 
extends the standard FLRW cosmology. This represents a fundamentally new 
approach to incorporating thermodynamic degrees of freedom into general 
relativity, moving beyond the traditional paradigm where quantities such 
as entropy, temperature, and mass appear only as parameters in the 
stress-energy tensor or equations of state. The AMMP framework provides 
the mathematical foundation for treating these physical quantities as 
coordinates on equal footing with space and time, potentially offering 
new insights into the deep connection between gravity and thermodynamics 
first suggested by black hole physics.

Our analysis demonstrates that the \textit{temporal-spatial-entropic metric} is not merely 
a mathematical curiosity but yields physically meaningful consequences. 
The Kretschmann scalar for the txyzS metric contains additional terms 
arising from the entropy dimension, including a constant entropy contribution 
$4(\beta S)^4$ and a coupling term $12H^2(\beta S)^2$ that directly 
connects cosmic expansion with entropy production. Numerical simulations 
reveal that these entropy-driven terms produce observable deviations from 
standard FLRW cosmology, particularly at late times and for larger entropy 
coordinates $S$. The ratio $K_{\text{txyzS}}/K_{\text{FLRW}}$ exceeds unity 
for all $S > 0$, with the deviation scaling as $S^2$ and $S^4$, suggesting 
that the entropy dimension amplifies spacetime curvature in a manner 
potentially detectable by cosmological observations.

This work establishes the AMMP framework as a viable and promising approach 
for generating novel metrics with physical consequences. The success of the 
entropy-dimension metric—its mathematical consistency, computational 
tractability, and production of observable deviations—validates the core 
idea of the AMMP framework: that thermodynamic quantities can be promoted 
to spacetime coordinates. Future work will extend this analysis to include 
temperature and mass as additional dimensions, explore the observational 
constraints on the $\beta$ parameter using CMB and large-scale structure 
data, and investigate the quantum aspects of the AMMP framework. The 
txyzS metric and its associated simulation code are freely available as 
part of the \href{https://github.com/lontelis/Einsteinpy_to_latex}{{\color{cyan}\texttt{EinsteinKPy2Tex}}}  
software package, enabling the community to 
build upon this foundation.

\section*{Acknowledgments}

The author thanks the EinsteinPy development team for their excellent library, and DeepSeek AI for assistance with code optimization and documentation.
PN is supported by NSFC Grant No. U2541210.

\section*{Data and Code Availability}

The software is open-source under the MIT license and available at:

\url{https://github.com/lontelis/EinsteinKPy2Tex}

\bibliographystyle{unsrt}
\bibliography{biblio}

\end{document}